% Methodology chapter for the final 30-city Feature-MAE experiment.
% Required packages in the parent manuscript:
% \usepackage{amsmath,amssymb}
% This file is a chapter fragment and can be included with:
% % Methodology chapter for the final 30-city Feature-MAE experiment.
% Required packages in the parent manuscript:
% \usepackage{amsmath,amssymb}
% This file is a chapter fragment and can be included with:
% % Methodology chapter for the final 30-city Feature-MAE experiment.
% Required packages in the parent manuscript:
% \usepackage{amsmath,amssymb}
% This file is a chapter fragment and can be included with:
% % Methodology chapter for the final 30-city Feature-MAE experiment.
% Required packages in the parent manuscript:
% \usepackage{amsmath,amssymb}
% This file is a chapter fragment and can be included with:
% \input{paper/methodology}

\section{Methodology}
\label{sec:methodology}

\subsection{Overview}
\label{sec:method_overview}

We developed a fully unsupervised framework for discovering multiscale urban
visual patterns from geotagged street-view imagery. The framework contains six
stages: (i) geographically broad and city-balanced image sampling; (ii) image
quality control; (iii) extraction of dense patch representations using a frozen
DINOv3 vision transformer; (iv) contextual representation learning using a
feature-level masked autoencoder (Feature-MAE); (v) spatial visualization and
hierarchical organization of the learned latent dimensions; and (vi) post-hoc
semantic interpretation using a compact vision--language model. Statistical
learning and clustering were completed before semantic labels were generated,
so language-model outputs could not affect the learned representation or the
cluster assignments.

The final analysis used Feature-MAE rather than the Top-$K$ sparse autoencoder
examined during model development. Both models used a latent width of 512, but
the preliminary sparse autoencoder retained the 32 largest activations for each
input, whereas the final Feature-MAE learned signed, contextual patch codes by
reconstructing masked DINOv3 features. Unless explicitly described as a pilot
or ablation, all analyses and figures reported in this study refer to the final
Feature-MAE model.

\subsection{Street-view image sample}
\label{sec:streetview_sample}

The study sample comprised street-view imagery from 30 cities spanning Africa,
Asia, Europe, North America, South America and Oceania: Amsterdam, Bangkok,
Bogot\'a, Buenos Aires, Cape Town, Dhaka, Hong Kong, Istanbul, Jakarta,
Johannesburg, Kuala Lumpur, Lagos, Lima, London, Los Angeles, Manila, Mexico
City, Moscow, Mumbai, New Delhi, New York City, Osaka, Paris, S\~ao Paulo,
Seoul, Singapore, Sydney, Taipei, Toronto and Vienna. Four nominal viewing
directions ($0^{\circ}$, $90^{\circ}$, $180^{\circ}$ and $270^{\circ}$) were
available for each panorama.

We first sampled 10,000 panoramas per city. From their directional views, we
then selected 12,800 images per city, yielding 384,000 images before quality
control. Selection used a deterministic, city-specific random seed derived from
the global seed of 42. Within each city, panorama order and heading order were
randomized and images were selected in a round-robin procedure: one randomly
ordered heading was selected from every panorama before a second heading was
selected from any panorama. This design retained all 10,000 unique panoramas in
each city while avoiding disproportionate sampling of multiple headings from a
small number of locations. It also maintained approximately balanced heading
counts within each city.

\subsection{Image quality control}
\label{sec:image_qc}

Image quality was assessed before representation learning to reduce the risk
that latent dimensions would encode acquisition defects rather than urban
content. Each image was resized to $256\times256$ pixels and described using 11
low-cost photometric and structural measurements: mean and standard deviation
of luminance, saturated-pixel fraction, dark-pixel fraction, Laplacian variance,
edge density, luminance entropy, dominant-intensity fraction, colourfulness,
mean red--green contrast and mean yellow--blue contrast. A histogram-based
gradient-boosting classifier, trained from conservative weak labels for clear
good- and bad-quality images, converted these measurements into an artefact
probability $p_{\mathrm{bad}}$. Images with $p_{\mathrm{bad}}>0.5$ were excluded.

We supplemented this classifier with explicit rules for unreadable, black,
blurred and tunnel-like frames. Let $\mu_Y$, $\sigma_Y$, $f_{Y<20}$ and
$V_{\nabla^2Y}$ denote mean luminance, luminance standard deviation, the
fraction of pixels with luminance below 20, and Laplacian variance, respectively.
An image was classified as black or uninformative if
$\mu_Y<14$, $f_{Y<20}>0.97$ or $\sigma_Y<6$. To detect spatially localized blur,
each image was divided into a $4\times4$ grid; an image was excluded when more
than 40\% of informative tiles had a tile-level Laplacian variance below 12, or
when the whole-image Laplacian variance was below 48. Tunnel-like views were
identified among non-black images using the conjunction
\begin{equation}
  \mu_Y < 100, \qquad
  \mu_{Y,\mathrm{top}} < 95, \qquad
  V_{\nabla^2Y} < 280,
  \label{eq:tunnel_gate}
\end{equation}
and either a red--blue mean difference greater than 4 or a highlight-glow
contrast greater than 42. The final rejection indicator was the union of the
learned artefact decision, the deterministic rules and image-read failure.

This procedure removed 5,182 of the 384,000 sampled images (1.35\%), leaving
378,818 images for model training and downstream analysis. Because rejection
reasons were not mutually exclusive, their counts were allowed to overlap. The
complete retained and excluded manifests were preserved to make the filtering
decision auditable.

\subsection{Dense DINOv3 feature extraction}
\label{sec:dinov3_features}

We used a frozen DINOv3 ViT-B/16 model pretrained on LVD-1689M as the visual
backbone \cite{simeoni2025dinov3}. Images were converted to RGB, resized to
$224\times224$ pixels with bicubic interpolation and normalized using the
ImageNet channel means $(0.485,0.456,0.406)$ and standard deviations
$(0.229,0.224,0.225)$. Given the patch size of 16 pixels, the image was
represented by a $14\times14$ grid of 196 patch tokens. We removed the class
token and all register tokens from the final hidden sequence and retained the
768-dimensional token associated with each spatial patch. Each token was
independently $\ell_2$-normalized:
\begin{equation}
  \mathbf{t}_{ip}
  = \frac{f_{\theta}(I_i)_p}
  {\lVert f_{\theta}(I_i)_p\rVert_2},
  \qquad
  \mathbf{T}_i\in\mathbb{R}^{196\times768},
  \label{eq:dino_tokens}
\end{equation}
where $i$ indexes images, $p$ indexes spatial patches, and $f_{\theta}$ denotes
the frozen DINOv3 encoder. Patch tensors were stored in half precision using
memory-mapped arrays; 75,264,000 tokens were extracted before quality filtering,
of which 74,248,328 belonged to retained images.

Feature extraction used batches of 128 images and bfloat16 automatic mixed
precision. Inference batches were recursively bisected following a CUDA
out-of-memory exception, providing a deterministic fallback for limited GPU
memory. Extraction was resumable at image level and required approximately
0.95~GiB of peak allocated GPU memory in the recorded run.

\subsection{Feature-level masked autoencoder}
\label{sec:feature_mae}

\subsubsection{Architecture}

We trained a feature-level masked autoencoder on the frozen DINOv3 patch tokens.
The design follows the asymmetric visible-token encoder and lightweight decoder
principle of masked autoencoding \cite{he2022maskedautoencoders}, but reconstructs
DINOv3 features rather than RGB pixels. For image $i$, an affine projection maps
each input token to a 512-dimensional latent space:
\begin{equation}
  \mathbf{z}^{(0)}_{ip}
  = \mathbf{W}_{\mathrm{in}}\mathbf{t}_{ip}
    + \mathbf{b}_{\mathrm{in}} + \boldsymbol{\pi}^{(e)}_p,
  \label{eq:mae_projection}
\end{equation}
where $\boldsymbol{\pi}^{(e)}_p$ is a fixed two-dimensional sine--cosine
positional embedding. For every training image, 75\% of the 196 positions were
randomly masked, leaving 49 visible tokens. Only visible tokens were processed
by the encoder:
\begin{equation}
  \mathbf{H}_{i,V}
  = \operatorname{LN}\!\left(
      E_{\phi}(\mathbf{Z}^{(0)}_{i,V})
    \right),
  \label{eq:mae_encoder}
\end{equation}
where $V$ is the visible-position set. The encoder consisted of four pre-norm
transformer blocks with a hidden width of 512, eight attention heads, a
feed-forward expansion ratio of four, GELU activations and no dropout.

Encoded visible tokens were projected from 512 to 256 dimensions. A learned
mask token was inserted at each missing position, the original spatial order was
restored, and a fixed 256-dimensional sine--cosine positional embedding was
added. A two-layer, eight-head transformer decoder with width 256 and a
fourfold feed-forward expansion produced reconstructed 768-dimensional patch
features:
\begin{equation}
  \widehat{\mathbf{T}}_i
  = \mathbf{W}_{\mathrm{out}}
    \operatorname{LN}\!\left[
      D_{\psi}\!\left(
        \operatorname{Restore}(
          \mathbf{W}_{ed}\mathbf{H}_{i,V},
          \mathbf{m})
        + \boldsymbol{\Pi}^{(d)}
      \right)
    \right]
    + \mathbf{b}_{\mathrm{out}},
  \label{eq:mae_decoder}
\end{equation}
where $\mathbf{m}$ is the learned mask token and
$\mathbf{W}_{ed}$ is the encoder-to-decoder projection. Linear layers were
initialized with Xavier-uniform weights, LayerNorm scales were initialized to
one, and the mask token was sampled from
$\mathcal{N}(0,0.02^2)$. The complete model contained 14,913,280 trainable
parameters.

\subsubsection{Training objective}

Because input patch tokens were normalized, reconstruction was optimized using
cosine distance only at masked positions. If $M_i$ is the masked-position set,
the loss for a batch $\mathcal{B}$ was
\begin{equation}
  \mathcal{L}_{\mathrm{mask}}
  = \frac{1}{\sum_{i\in\mathcal{B}}|M_i|}
    \sum_{i\in\mathcal{B}}
    \sum_{p\in M_i}
    \left[
      1-
      \frac{\mathbf{t}_{ip}^{\top}\widehat{\mathbf{t}}_{ip}}
      {\lVert\mathbf{t}_{ip}\rVert_2
       \lVert\widehat{\mathbf{t}}_{ip}\rVert_2}
    \right].
  \label{eq:masked_cosine_loss}
\end{equation}
No semantic supervision, city label, Top-$K$ operator, non-negativity penalty or
clustering loss was used to train the final Feature-MAE.

\subsubsection{City-balanced optimization}

We reserved 128 clean images per city for validation, producing a validation
set of 3,840 images and a distinct-image training pool of 374,978 images. Each
training batch contained the same number of images from every city. Cities with
fewer retained images were reshuffled and resampled within an epoch until they
matched the largest city-specific training set, preventing cities with more
retained images from dominating optimization.

The model was trained for ten epochs using AdamW with
$\beta_1=0.9$, $\beta_2=0.95$, an initial learning rate of
$1.5\times10^{-4}$ and weight decay 0.05. A linear warm-up covered the first
5\% of optimizer steps, followed by cosine decay to zero. Gradients were clipped
to a global norm of one. Training used bfloat16 automatic mixed precision and
TensorFloat-32 matrix multiplication. The city-balanced batch contained eight
images per city (240 images) during epochs 1--4 and 16 images per city (480
images) during epochs 5--10; the latter was processed as two microbatches of 240
images. The checkpoint with the lowest validation loss was retained. The selected
checkpoint was obtained at epoch 10, with masked cosine losses of 0.10356 and
0.10334 on the training and validation sets, respectively. Peak allocated GPU
memory was 3.27~GiB.

\subsection{Latent activation maps and representative images}
\label{sec:latent_heatmaps}

At inference, masking and the decoder were removed. All 196 DINOv3 patch tokens
were projected and processed by the trained encoder:
\begin{equation}
  \mathbf{A}_i
  = \operatorname{LN}\!\left[
      E_{\phi}(\mathbf{W}_{\mathrm{in}}\mathbf{T}_i
      +\mathbf{b}_{\mathrm{in}}+\boldsymbol{\Pi}^{(e)})
    \right]
  \in \mathbb{R}^{196\times512}.
  \label{eq:full_encoding}
\end{equation}
Reshaping $\mathbf{A}_i$ yielded a
$14\times14\times512$ spatial activation tensor. The activations were signed;
therefore, dimension $d$ was visualized using its positive evidence
$\max(A_{ipd},0)$ rather than the absolute activation. For each dimension, we
defined the image-level response
\begin{equation}
  s_{id}=\max_{p\in\{1,\ldots,196\}} A_{ipd}.
  \label{eq:image_response}
\end{equation}
All 378,818 retained images were scanned. The 128 highest-scoring candidates
were retained provisionally, after which duplicate panoramas were removed and
the ten highest-scoring unique panoramas were selected. This ensured that a
dimension was not interpreted from multiple headings of the same location.

For visualization, positive activations were divided by the maximum response of
the selected examples, clipped to $[0,1]$, mapped to a perceptually ordered
blue--cyan--yellow--red colour scale, and resized to the original image size
using bicubic interpolation. Transparency increased monotonically with
activation so that near-zero patches remained visually unchanged. The resulting
heatmap therefore indicates where a dimension responds, whereas the original
image remains the source for interpreting scene content.

\subsection{Multiview similarity between latent dimensions}
\label{sec:multiview_similarity}

A single parameter-space similarity is insufficient to characterize a contextual
latent dimension. We therefore represented each of the 512 dimensions through
five complementary views: encoder direction, linearized decoder direction,
patch-position profile, city profile and activation-context profile.

First, the encoder direction of dimension $d$ was the $d$th row of the input
projection matrix,
\begin{equation}
  \mathbf{e}_d = (\mathbf{W}_{\mathrm{in}})_{d,:}
  \in\mathbb{R}^{768}.
  \label{eq:encoder_direction}
\end{equation}
Second, we linearized the path from encoder latents to reconstructed DINOv3
features, obtaining
\begin{equation}
  \mathbf{g}_d
  = \left[
      (\mathbf{W}_{\mathrm{out}}\mathbf{W}_{ed})^{\top}
    \right]_{d,:}
  \in\mathbb{R}^{768}.
  \label{eq:decoder_direction}
\end{equation}

For the remaining views, each patch was assigned to its largest latent
activation,
\begin{equation}
  w_{ip}=\underset{d\in\{1,\ldots,512\}}{\arg\max}\; A_{ipd}.
  \label{eq:winning_dimension}
\end{equation}
The patch-position profile
$\mathbf{p}_d\in\mathbb{R}^{196}$ counted how often dimension $d$ won at each
grid position, while the city profile
$\mathbf{c}_d\in\mathbb{R}^{30}$ counted winning patches by city:
\begin{align}
  p_{dq} &= \sum_i \mathbb{I}(w_{iq}=d), \\
  c_{du} &= \sum_{i:\,\operatorname{city}(i)=u}\sum_p
            \mathbb{I}(w_{ip}=d).
  \label{eq:position_city_profiles}
\end{align}

Finally, we constructed an image-level activation-context view. Let
$B_{id}=1$ if dimension $d$ won at least one patch in image $i$, and zero
otherwise. For $N$ images, define the co-occurrence count
$n_{de}=\sum_iB_{id}B_{ie}$ and support $n_d=\sum_iB_{id}$. We used a smoothed
positive pointwise mutual information matrix
\begin{equation}
  Q_{de}=\max\left\{
    0,
    \log\frac{n_{de}+0.5}{n_dn_e/N+0.5}
  \right\}, \qquad Q_{dd}=0,
  \label{eq:ppmi}
\end{equation}
and treated row $\mathbf{q}_d$ as the contextual profile of dimension $d$.

For each view $v\in\{\mathbf{e},\mathbf{g},\mathbf{p},\mathbf{c},\mathbf{q}\}$,
we computed a $512\times512$ cosine-similarity matrix:
\begin{equation}
  S^{(v)}_{de}
  =\frac{\mathbf{v}_d^{\top}\mathbf{v}_e}
  {\lVert\mathbf{v}_d\rVert_2\lVert\mathbf{v}_e\rVert_2}.
  \label{eq:view_cosine}
\end{equation}

\subsection{Rank fusion and hierarchical clustering}
\label{sec:hierarchical_clustering}

The five similarity views have different numerical distributions, so we fused
neighbour ranks rather than raw similarity values. For every dimension $d$ and
view $v$, the 30 nearest non-self neighbours were ordered by decreasing
similarity. A neighbour at zero-indexed rank $r\in\{0,\ldots,29\}$ received
weight $1-r/30$. The directed fused affinity was
\begin{equation}
  \widetilde{F}_{de}
  =\frac{1}{5}\sum_v
  \mathbb{I}\!\left[e\in\mathcal{N}^{(v)}_{30}(d)\right]
  \left(1-\frac{r^{(v)}_{de}}{30}\right).
  \label{eq:rank_fusion}
\end{equation}
We then symmetrized and rescaled the matrix:
\begin{equation}
  F=\frac{\widetilde{F}+\widetilde{F}^{\top}}{2},
  \qquad F_{dd}=0, \qquad
  F\leftarrow\frac{F}{\max_{d,e}F_{de}}.
  \label{eq:fused_affinity}
\end{equation}

To preserve nonlinear neighbourhood structure, the fused affinity was projected
into 32 dimensions using the eigenvectors of the symmetric normalized graph
Laplacian \cite{ng2001spectral}. The first eigenvector was retained, consistent
with the spectral-clustering embedding used in the implementation. Ward's
minimum-variance agglomerative method was then applied to Euclidean distances in
the spectral embedding \cite{ward1963hierarchical}. One dendrogram was
constructed and cut at two resolutions: 32 coarse clusters and 64 fine clusters.
Because both partitions were cuts of the same tree, each fine cluster had
exactly one coarse parent by construction.

For cluster $C$, the representative or medoid dimension was defined using the
fused affinity rather than the spectral coordinates:
\begin{equation}
  d_C^{*}=\underset{d\in C}{\arg\max}
  \sum_{e\in C}F_{de}.
  \label{eq:cluster_medoid}
\end{equation}
Cluster identities were assigned only after the tree had been fixed. Canonical
identifiers C000--C031 and F000--F063 were ordered by the first latent dimension
appearing in each cluster and carry no semantic meaning.

\subsection{Cluster validity and split-sample stability}
\label{sec:cluster_validation}

We characterized both hierarchy levels using silhouette score,
Calinski--Harabasz index, Davies--Bouldin index and normalized cluster-size
entropy. If $n_k$ denotes the size of cluster $k$, $p_k=n_k/512$ and $K$ is the
number of clusters, normalized size entropy was
\begin{equation}
  H_{\mathrm{size}}
  =-\frac{\sum_{k=1}^{K}p_k\log p_k}{\log K}.
  \label{eq:size_entropy}
\end{equation}

Stability was evaluated using two deterministic, interleaved halves of the image
sample. Parameter-space views (encoder and decoder directions) were held fixed,
whereas position, city and PPMI-context profiles were independently recomputed
for each half. The complete rank-fusion, spectral-embedding and Ward-clustering
pipeline was repeated for each half using seeds 43 and 44. For each dimension,
we calculated the Jaccard overlap between its full-data cluster and its
corresponding half-sample cluster. Global agreement was additionally summarized
with the adjusted Rand index and normalized mutual information at both the
32- and 64-cluster resolutions. These stability measures quantify sensitivity to
the image sample and were not used to choose or modify semantic labels.

\subsection{Post-hoc semantic interpretation}
\label{sec:semantic_interpretation}

We used the locally deployed Qwen3-VL-2B-Instruct model
\cite{bai2025qwen3vl} to attach human-readable descriptions to the statistical
clusters. Evidence sheets were resized to a maximum side length of 1,200 pixels,
and generation was limited to 240 new tokens. Interpretation was performed after
Feature-MAE training and after the
32/64 hierarchy had been frozen. The language model therefore had no access to
the optimization objective and could not influence dimension similarity or
cluster membership.

For each of the 64 fine and 32 coarse clusters, we selected the cluster medoid
from Eq.~\ref{eq:cluster_medoid}. We then retrieved that dimension's ten
highest-scoring images from unique panoramas. Each item in the evidence sheet
contained (i) the original street-view image, (ii) the activation overlay and
(iii) a crop centred on the maximally activated patch. The model was instructed
to identify visual evidence that recurred across the ten examples and to return
a strict JSON object containing a concise Chinese name, a short English label,
a semantic type, a summary, three visual cues, a tentative urban interpretation,
cross-image consistency, artefact probability and confidence. Semantic type was
restricted to object, material, colour/light, texture, geometry, spatial layout,
environment, artefact, mixed or unclear.

Fine clusters were interpreted first. When a coarse cluster was interpreted,
the already generated descriptions of its fine children were supplied together
with the coarse medoid evidence, allowing the model to formulate a broader
parent concept without changing the statistical hierarchy. Outputs were parsed
as JSON and automatically retried when parsing failed. A second generation was
requested when a name contained prohibited generic terms, confidence was
mechanically set near zero, or the reported visual cues were duplicated. Exactly
ten images were used for each of the 96 categories; interpreting every training
image was neither necessary nor performed.

\subsection{Implementation and reproducibility}
\label{sec:implementation}

All stochastic components used a global seed of 42 unless a split-sample seed
was explicitly specified. DINOv3 and Qwen3-VL weights were loaded locally, and
the DINOv3 backbone remained frozen throughout. Large arrays were stored as
half-precision memory maps, and image selection, feature extraction, training,
ranking and clustering wrote resumable intermediate states. GPU-intensive
PyTorch stages and CPU-intensive SciPy/scikit-learn stages were executed in
separate processes to avoid conflicts between OpenMP runtimes. The statistical
hierarchy, category membership table, linkage matrix, fused similarity matrix,
spectral embedding, heatmap metadata, semantic prompts and raw language-model
responses were retained to support independent auditing and figure
reconstruction.


\section{Methodology}
\label{sec:methodology}

\subsection{Overview}
\label{sec:method_overview}

We developed a fully unsupervised framework for discovering multiscale urban
visual patterns from geotagged street-view imagery. The framework contains six
stages: (i) geographically broad and city-balanced image sampling; (ii) image
quality control; (iii) extraction of dense patch representations using a frozen
DINOv3 vision transformer; (iv) contextual representation learning using a
feature-level masked autoencoder (Feature-MAE); (v) spatial visualization and
hierarchical organization of the learned latent dimensions; and (vi) post-hoc
semantic interpretation using a compact vision--language model. Statistical
learning and clustering were completed before semantic labels were generated,
so language-model outputs could not affect the learned representation or the
cluster assignments.

The final analysis used Feature-MAE rather than the Top-$K$ sparse autoencoder
examined during model development. Both models used a latent width of 512, but
the preliminary sparse autoencoder retained the 32 largest activations for each
input, whereas the final Feature-MAE learned signed, contextual patch codes by
reconstructing masked DINOv3 features. Unless explicitly described as a pilot
or ablation, all analyses and figures reported in this study refer to the final
Feature-MAE model.

\subsection{Street-view image sample}
\label{sec:streetview_sample}

The study sample comprised street-view imagery from 30 cities spanning Africa,
Asia, Europe, North America, South America and Oceania: Amsterdam, Bangkok,
Bogot\'a, Buenos Aires, Cape Town, Dhaka, Hong Kong, Istanbul, Jakarta,
Johannesburg, Kuala Lumpur, Lagos, Lima, London, Los Angeles, Manila, Mexico
City, Moscow, Mumbai, New Delhi, New York City, Osaka, Paris, S\~ao Paulo,
Seoul, Singapore, Sydney, Taipei, Toronto and Vienna. Four nominal viewing
directions ($0^{\circ}$, $90^{\circ}$, $180^{\circ}$ and $270^{\circ}$) were
available for each panorama.

We first sampled 10,000 panoramas per city. From their directional views, we
then selected 12,800 images per city, yielding 384,000 images before quality
control. Selection used a deterministic, city-specific random seed derived from
the global seed of 42. Within each city, panorama order and heading order were
randomized and images were selected in a round-robin procedure: one randomly
ordered heading was selected from every panorama before a second heading was
selected from any panorama. This design retained all 10,000 unique panoramas in
each city while avoiding disproportionate sampling of multiple headings from a
small number of locations. It also maintained approximately balanced heading
counts within each city.

\subsection{Image quality control}
\label{sec:image_qc}

Image quality was assessed before representation learning to reduce the risk
that latent dimensions would encode acquisition defects rather than urban
content. Each image was resized to $256\times256$ pixels and described using 11
low-cost photometric and structural measurements: mean and standard deviation
of luminance, saturated-pixel fraction, dark-pixel fraction, Laplacian variance,
edge density, luminance entropy, dominant-intensity fraction, colourfulness,
mean red--green contrast and mean yellow--blue contrast. A histogram-based
gradient-boosting classifier, trained from conservative weak labels for clear
good- and bad-quality images, converted these measurements into an artefact
probability $p_{\mathrm{bad}}$. Images with $p_{\mathrm{bad}}>0.5$ were excluded.

We supplemented this classifier with explicit rules for unreadable, black,
blurred and tunnel-like frames. Let $\mu_Y$, $\sigma_Y$, $f_{Y<20}$ and
$V_{\nabla^2Y}$ denote mean luminance, luminance standard deviation, the
fraction of pixels with luminance below 20, and Laplacian variance, respectively.
An image was classified as black or uninformative if
$\mu_Y<14$, $f_{Y<20}>0.97$ or $\sigma_Y<6$. To detect spatially localized blur,
each image was divided into a $4\times4$ grid; an image was excluded when more
than 40\% of informative tiles had a tile-level Laplacian variance below 12, or
when the whole-image Laplacian variance was below 48. Tunnel-like views were
identified among non-black images using the conjunction
\begin{equation}
  \mu_Y < 100, \qquad
  \mu_{Y,\mathrm{top}} < 95, \qquad
  V_{\nabla^2Y} < 280,
  \label{eq:tunnel_gate}
\end{equation}
and either a red--blue mean difference greater than 4 or a highlight-glow
contrast greater than 42. The final rejection indicator was the union of the
learned artefact decision, the deterministic rules and image-read failure.

This procedure removed 5,182 of the 384,000 sampled images (1.35\%), leaving
378,818 images for model training and downstream analysis. Because rejection
reasons were not mutually exclusive, their counts were allowed to overlap. The
complete retained and excluded manifests were preserved to make the filtering
decision auditable.

\subsection{Dense DINOv3 feature extraction}
\label{sec:dinov3_features}

We used a frozen DINOv3 ViT-B/16 model pretrained on LVD-1689M as the visual
backbone \cite{simeoni2025dinov3}. Images were converted to RGB, resized to
$224\times224$ pixels with bicubic interpolation and normalized using the
ImageNet channel means $(0.485,0.456,0.406)$ and standard deviations
$(0.229,0.224,0.225)$. Given the patch size of 16 pixels, the image was
represented by a $14\times14$ grid of 196 patch tokens. We removed the class
token and all register tokens from the final hidden sequence and retained the
768-dimensional token associated with each spatial patch. Each token was
independently $\ell_2$-normalized:
\begin{equation}
  \mathbf{t}_{ip}
  = \frac{f_{\theta}(I_i)_p}
  {\lVert f_{\theta}(I_i)_p\rVert_2},
  \qquad
  \mathbf{T}_i\in\mathbb{R}^{196\times768},
  \label{eq:dino_tokens}
\end{equation}
where $i$ indexes images, $p$ indexes spatial patches, and $f_{\theta}$ denotes
the frozen DINOv3 encoder. Patch tensors were stored in half precision using
memory-mapped arrays; 75,264,000 tokens were extracted before quality filtering,
of which 74,248,328 belonged to retained images.

Feature extraction used batches of 128 images and bfloat16 automatic mixed
precision. Inference batches were recursively bisected following a CUDA
out-of-memory exception, providing a deterministic fallback for limited GPU
memory. Extraction was resumable at image level and required approximately
0.95~GiB of peak allocated GPU memory in the recorded run.

\subsection{Feature-level masked autoencoder}
\label{sec:feature_mae}

\subsubsection{Architecture}

We trained a feature-level masked autoencoder on the frozen DINOv3 patch tokens.
The design follows the asymmetric visible-token encoder and lightweight decoder
principle of masked autoencoding \cite{he2022maskedautoencoders}, but reconstructs
DINOv3 features rather than RGB pixels. For image $i$, an affine projection maps
each input token to a 512-dimensional latent space:
\begin{equation}
  \mathbf{z}^{(0)}_{ip}
  = \mathbf{W}_{\mathrm{in}}\mathbf{t}_{ip}
    + \mathbf{b}_{\mathrm{in}} + \boldsymbol{\pi}^{(e)}_p,
  \label{eq:mae_projection}
\end{equation}
where $\boldsymbol{\pi}^{(e)}_p$ is a fixed two-dimensional sine--cosine
positional embedding. For every training image, 75\% of the 196 positions were
randomly masked, leaving 49 visible tokens. Only visible tokens were processed
by the encoder:
\begin{equation}
  \mathbf{H}_{i,V}
  = \operatorname{LN}\!\left(
      E_{\phi}(\mathbf{Z}^{(0)}_{i,V})
    \right),
  \label{eq:mae_encoder}
\end{equation}
where $V$ is the visible-position set. The encoder consisted of four pre-norm
transformer blocks with a hidden width of 512, eight attention heads, a
feed-forward expansion ratio of four, GELU activations and no dropout.

Encoded visible tokens were projected from 512 to 256 dimensions. A learned
mask token was inserted at each missing position, the original spatial order was
restored, and a fixed 256-dimensional sine--cosine positional embedding was
added. A two-layer, eight-head transformer decoder with width 256 and a
fourfold feed-forward expansion produced reconstructed 768-dimensional patch
features:
\begin{equation}
  \widehat{\mathbf{T}}_i
  = \mathbf{W}_{\mathrm{out}}
    \operatorname{LN}\!\left[
      D_{\psi}\!\left(
        \operatorname{Restore}(
          \mathbf{W}_{ed}\mathbf{H}_{i,V},
          \mathbf{m})
        + \boldsymbol{\Pi}^{(d)}
      \right)
    \right]
    + \mathbf{b}_{\mathrm{out}},
  \label{eq:mae_decoder}
\end{equation}
where $\mathbf{m}$ is the learned mask token and
$\mathbf{W}_{ed}$ is the encoder-to-decoder projection. Linear layers were
initialized with Xavier-uniform weights, LayerNorm scales were initialized to
one, and the mask token was sampled from
$\mathcal{N}(0,0.02^2)$. The complete model contained 14,913,280 trainable
parameters.

\subsubsection{Training objective}

Because input patch tokens were normalized, reconstruction was optimized using
cosine distance only at masked positions. If $M_i$ is the masked-position set,
the loss for a batch $\mathcal{B}$ was
\begin{equation}
  \mathcal{L}_{\mathrm{mask}}
  = \frac{1}{\sum_{i\in\mathcal{B}}|M_i|}
    \sum_{i\in\mathcal{B}}
    \sum_{p\in M_i}
    \left[
      1-
      \frac{\mathbf{t}_{ip}^{\top}\widehat{\mathbf{t}}_{ip}}
      {\lVert\mathbf{t}_{ip}\rVert_2
       \lVert\widehat{\mathbf{t}}_{ip}\rVert_2}
    \right].
  \label{eq:masked_cosine_loss}
\end{equation}
No semantic supervision, city label, Top-$K$ operator, non-negativity penalty or
clustering loss was used to train the final Feature-MAE.

\subsubsection{City-balanced optimization}

We reserved 128 clean images per city for validation, producing a validation
set of 3,840 images and a distinct-image training pool of 374,978 images. Each
training batch contained the same number of images from every city. Cities with
fewer retained images were reshuffled and resampled within an epoch until they
matched the largest city-specific training set, preventing cities with more
retained images from dominating optimization.

The model was trained for ten epochs using AdamW with
$\beta_1=0.9$, $\beta_2=0.95$, an initial learning rate of
$1.5\times10^{-4}$ and weight decay 0.05. A linear warm-up covered the first
5\% of optimizer steps, followed by cosine decay to zero. Gradients were clipped
to a global norm of one. Training used bfloat16 automatic mixed precision and
TensorFloat-32 matrix multiplication. The city-balanced batch contained eight
images per city (240 images) during epochs 1--4 and 16 images per city (480
images) during epochs 5--10; the latter was processed as two microbatches of 240
images. The checkpoint with the lowest validation loss was retained. The selected
checkpoint was obtained at epoch 10, with masked cosine losses of 0.10356 and
0.10334 on the training and validation sets, respectively. Peak allocated GPU
memory was 3.27~GiB.

\subsection{Latent activation maps and representative images}
\label{sec:latent_heatmaps}

At inference, masking and the decoder were removed. All 196 DINOv3 patch tokens
were projected and processed by the trained encoder:
\begin{equation}
  \mathbf{A}_i
  = \operatorname{LN}\!\left[
      E_{\phi}(\mathbf{W}_{\mathrm{in}}\mathbf{T}_i
      +\mathbf{b}_{\mathrm{in}}+\boldsymbol{\Pi}^{(e)})
    \right]
  \in \mathbb{R}^{196\times512}.
  \label{eq:full_encoding}
\end{equation}
Reshaping $\mathbf{A}_i$ yielded a
$14\times14\times512$ spatial activation tensor. The activations were signed;
therefore, dimension $d$ was visualized using its positive evidence
$\max(A_{ipd},0)$ rather than the absolute activation. For each dimension, we
defined the image-level response
\begin{equation}
  s_{id}=\max_{p\in\{1,\ldots,196\}} A_{ipd}.
  \label{eq:image_response}
\end{equation}
All 378,818 retained images were scanned. The 128 highest-scoring candidates
were retained provisionally, after which duplicate panoramas were removed and
the ten highest-scoring unique panoramas were selected. This ensured that a
dimension was not interpreted from multiple headings of the same location.

For visualization, positive activations were divided by the maximum response of
the selected examples, clipped to $[0,1]$, mapped to a perceptually ordered
blue--cyan--yellow--red colour scale, and resized to the original image size
using bicubic interpolation. Transparency increased monotonically with
activation so that near-zero patches remained visually unchanged. The resulting
heatmap therefore indicates where a dimension responds, whereas the original
image remains the source for interpreting scene content.

\subsection{Multiview similarity between latent dimensions}
\label{sec:multiview_similarity}

A single parameter-space similarity is insufficient to characterize a contextual
latent dimension. We therefore represented each of the 512 dimensions through
five complementary views: encoder direction, linearized decoder direction,
patch-position profile, city profile and activation-context profile.

First, the encoder direction of dimension $d$ was the $d$th row of the input
projection matrix,
\begin{equation}
  \mathbf{e}_d = (\mathbf{W}_{\mathrm{in}})_{d,:}
  \in\mathbb{R}^{768}.
  \label{eq:encoder_direction}
\end{equation}
Second, we linearized the path from encoder latents to reconstructed DINOv3
features, obtaining
\begin{equation}
  \mathbf{g}_d
  = \left[
      (\mathbf{W}_{\mathrm{out}}\mathbf{W}_{ed})^{\top}
    \right]_{d,:}
  \in\mathbb{R}^{768}.
  \label{eq:decoder_direction}
\end{equation}

For the remaining views, each patch was assigned to its largest latent
activation,
\begin{equation}
  w_{ip}=\underset{d\in\{1,\ldots,512\}}{\arg\max}\; A_{ipd}.
  \label{eq:winning_dimension}
\end{equation}
The patch-position profile
$\mathbf{p}_d\in\mathbb{R}^{196}$ counted how often dimension $d$ won at each
grid position, while the city profile
$\mathbf{c}_d\in\mathbb{R}^{30}$ counted winning patches by city:
\begin{align}
  p_{dq} &= \sum_i \mathbb{I}(w_{iq}=d), \\
  c_{du} &= \sum_{i:\,\operatorname{city}(i)=u}\sum_p
            \mathbb{I}(w_{ip}=d).
  \label{eq:position_city_profiles}
\end{align}

Finally, we constructed an image-level activation-context view. Let
$B_{id}=1$ if dimension $d$ won at least one patch in image $i$, and zero
otherwise. For $N$ images, define the co-occurrence count
$n_{de}=\sum_iB_{id}B_{ie}$ and support $n_d=\sum_iB_{id}$. We used a smoothed
positive pointwise mutual information matrix
\begin{equation}
  Q_{de}=\max\left\{
    0,
    \log\frac{n_{de}+0.5}{n_dn_e/N+0.5}
  \right\}, \qquad Q_{dd}=0,
  \label{eq:ppmi}
\end{equation}
and treated row $\mathbf{q}_d$ as the contextual profile of dimension $d$.

For each view $v\in\{\mathbf{e},\mathbf{g},\mathbf{p},\mathbf{c},\mathbf{q}\}$,
we computed a $512\times512$ cosine-similarity matrix:
\begin{equation}
  S^{(v)}_{de}
  =\frac{\mathbf{v}_d^{\top}\mathbf{v}_e}
  {\lVert\mathbf{v}_d\rVert_2\lVert\mathbf{v}_e\rVert_2}.
  \label{eq:view_cosine}
\end{equation}

\subsection{Rank fusion and hierarchical clustering}
\label{sec:hierarchical_clustering}

The five similarity views have different numerical distributions, so we fused
neighbour ranks rather than raw similarity values. For every dimension $d$ and
view $v$, the 30 nearest non-self neighbours were ordered by decreasing
similarity. A neighbour at zero-indexed rank $r\in\{0,\ldots,29\}$ received
weight $1-r/30$. The directed fused affinity was
\begin{equation}
  \widetilde{F}_{de}
  =\frac{1}{5}\sum_v
  \mathbb{I}\!\left[e\in\mathcal{N}^{(v)}_{30}(d)\right]
  \left(1-\frac{r^{(v)}_{de}}{30}\right).
  \label{eq:rank_fusion}
\end{equation}
We then symmetrized and rescaled the matrix:
\begin{equation}
  F=\frac{\widetilde{F}+\widetilde{F}^{\top}}{2},
  \qquad F_{dd}=0, \qquad
  F\leftarrow\frac{F}{\max_{d,e}F_{de}}.
  \label{eq:fused_affinity}
\end{equation}

To preserve nonlinear neighbourhood structure, the fused affinity was projected
into 32 dimensions using the eigenvectors of the symmetric normalized graph
Laplacian \cite{ng2001spectral}. The first eigenvector was retained, consistent
with the spectral-clustering embedding used in the implementation. Ward's
minimum-variance agglomerative method was then applied to Euclidean distances in
the spectral embedding \cite{ward1963hierarchical}. One dendrogram was
constructed and cut at two resolutions: 32 coarse clusters and 64 fine clusters.
Because both partitions were cuts of the same tree, each fine cluster had
exactly one coarse parent by construction.

For cluster $C$, the representative or medoid dimension was defined using the
fused affinity rather than the spectral coordinates:
\begin{equation}
  d_C^{*}=\underset{d\in C}{\arg\max}
  \sum_{e\in C}F_{de}.
  \label{eq:cluster_medoid}
\end{equation}
Cluster identities were assigned only after the tree had been fixed. Canonical
identifiers C000--C031 and F000--F063 were ordered by the first latent dimension
appearing in each cluster and carry no semantic meaning.

\subsection{Cluster validity and split-sample stability}
\label{sec:cluster_validation}

We characterized both hierarchy levels using silhouette score,
Calinski--Harabasz index, Davies--Bouldin index and normalized cluster-size
entropy. If $n_k$ denotes the size of cluster $k$, $p_k=n_k/512$ and $K$ is the
number of clusters, normalized size entropy was
\begin{equation}
  H_{\mathrm{size}}
  =-\frac{\sum_{k=1}^{K}p_k\log p_k}{\log K}.
  \label{eq:size_entropy}
\end{equation}

Stability was evaluated using two deterministic, interleaved halves of the image
sample. Parameter-space views (encoder and decoder directions) were held fixed,
whereas position, city and PPMI-context profiles were independently recomputed
for each half. The complete rank-fusion, spectral-embedding and Ward-clustering
pipeline was repeated for each half using seeds 43 and 44. For each dimension,
we calculated the Jaccard overlap between its full-data cluster and its
corresponding half-sample cluster. Global agreement was additionally summarized
with the adjusted Rand index and normalized mutual information at both the
32- and 64-cluster resolutions. These stability measures quantify sensitivity to
the image sample and were not used to choose or modify semantic labels.

\subsection{Post-hoc semantic interpretation}
\label{sec:semantic_interpretation}

We used the locally deployed Qwen3-VL-2B-Instruct model
\cite{bai2025qwen3vl} to attach human-readable descriptions to the statistical
clusters. Evidence sheets were resized to a maximum side length of 1,200 pixels,
and generation was limited to 240 new tokens. Interpretation was performed after
Feature-MAE training and after the
32/64 hierarchy had been frozen. The language model therefore had no access to
the optimization objective and could not influence dimension similarity or
cluster membership.

For each of the 64 fine and 32 coarse clusters, we selected the cluster medoid
from Eq.~\ref{eq:cluster_medoid}. We then retrieved that dimension's ten
highest-scoring images from unique panoramas. Each item in the evidence sheet
contained (i) the original street-view image, (ii) the activation overlay and
(iii) a crop centred on the maximally activated patch. The model was instructed
to identify visual evidence that recurred across the ten examples and to return
a strict JSON object containing a concise Chinese name, a short English label,
a semantic type, a summary, three visual cues, a tentative urban interpretation,
cross-image consistency, artefact probability and confidence. Semantic type was
restricted to object, material, colour/light, texture, geometry, spatial layout,
environment, artefact, mixed or unclear.

Fine clusters were interpreted first. When a coarse cluster was interpreted,
the already generated descriptions of its fine children were supplied together
with the coarse medoid evidence, allowing the model to formulate a broader
parent concept without changing the statistical hierarchy. Outputs were parsed
as JSON and automatically retried when parsing failed. A second generation was
requested when a name contained prohibited generic terms, confidence was
mechanically set near zero, or the reported visual cues were duplicated. Exactly
ten images were used for each of the 96 categories; interpreting every training
image was neither necessary nor performed.

\subsection{Implementation and reproducibility}
\label{sec:implementation}

All stochastic components used a global seed of 42 unless a split-sample seed
was explicitly specified. DINOv3 and Qwen3-VL weights were loaded locally, and
the DINOv3 backbone remained frozen throughout. Large arrays were stored as
half-precision memory maps, and image selection, feature extraction, training,
ranking and clustering wrote resumable intermediate states. GPU-intensive
PyTorch stages and CPU-intensive SciPy/scikit-learn stages were executed in
separate processes to avoid conflicts between OpenMP runtimes. The statistical
hierarchy, category membership table, linkage matrix, fused similarity matrix,
spectral embedding, heatmap metadata, semantic prompts and raw language-model
responses were retained to support independent auditing and figure
reconstruction.


\section{Methodology}
\label{sec:methodology}

\subsection{Overview}
\label{sec:method_overview}

We developed a fully unsupervised framework for discovering multiscale urban
visual patterns from geotagged street-view imagery. The framework contains six
stages: (i) geographically broad and city-balanced image sampling; (ii) image
quality control; (iii) extraction of dense patch representations using a frozen
DINOv3 vision transformer; (iv) contextual representation learning using a
feature-level masked autoencoder (Feature-MAE); (v) spatial visualization and
hierarchical organization of the learned latent dimensions; and (vi) post-hoc
semantic interpretation using a compact vision--language model. Statistical
learning and clustering were completed before semantic labels were generated,
so language-model outputs could not affect the learned representation or the
cluster assignments.

The final analysis used Feature-MAE rather than the Top-$K$ sparse autoencoder
examined during model development. Both models used a latent width of 512, but
the preliminary sparse autoencoder retained the 32 largest activations for each
input, whereas the final Feature-MAE learned signed, contextual patch codes by
reconstructing masked DINOv3 features. Unless explicitly described as a pilot
or ablation, all analyses and figures reported in this study refer to the final
Feature-MAE model.

\subsection{Street-view image sample}
\label{sec:streetview_sample}

The study sample comprised street-view imagery from 30 cities spanning Africa,
Asia, Europe, North America, South America and Oceania: Amsterdam, Bangkok,
Bogot\'a, Buenos Aires, Cape Town, Dhaka, Hong Kong, Istanbul, Jakarta,
Johannesburg, Kuala Lumpur, Lagos, Lima, London, Los Angeles, Manila, Mexico
City, Moscow, Mumbai, New Delhi, New York City, Osaka, Paris, S\~ao Paulo,
Seoul, Singapore, Sydney, Taipei, Toronto and Vienna. Four nominal viewing
directions ($0^{\circ}$, $90^{\circ}$, $180^{\circ}$ and $270^{\circ}$) were
available for each panorama.

We first sampled 10,000 panoramas per city. From their directional views, we
then selected 12,800 images per city, yielding 384,000 images before quality
control. Selection used a deterministic, city-specific random seed derived from
the global seed of 42. Within each city, panorama order and heading order were
randomized and images were selected in a round-robin procedure: one randomly
ordered heading was selected from every panorama before a second heading was
selected from any panorama. This design retained all 10,000 unique panoramas in
each city while avoiding disproportionate sampling of multiple headings from a
small number of locations. It also maintained approximately balanced heading
counts within each city.

\subsection{Image quality control}
\label{sec:image_qc}

Image quality was assessed before representation learning to reduce the risk
that latent dimensions would encode acquisition defects rather than urban
content. Each image was resized to $256\times256$ pixels and described using 11
low-cost photometric and structural measurements: mean and standard deviation
of luminance, saturated-pixel fraction, dark-pixel fraction, Laplacian variance,
edge density, luminance entropy, dominant-intensity fraction, colourfulness,
mean red--green contrast and mean yellow--blue contrast. A histogram-based
gradient-boosting classifier, trained from conservative weak labels for clear
good- and bad-quality images, converted these measurements into an artefact
probability $p_{\mathrm{bad}}$. Images with $p_{\mathrm{bad}}>0.5$ were excluded.

We supplemented this classifier with explicit rules for unreadable, black,
blurred and tunnel-like frames. Let $\mu_Y$, $\sigma_Y$, $f_{Y<20}$ and
$V_{\nabla^2Y}$ denote mean luminance, luminance standard deviation, the
fraction of pixels with luminance below 20, and Laplacian variance, respectively.
An image was classified as black or uninformative if
$\mu_Y<14$, $f_{Y<20}>0.97$ or $\sigma_Y<6$. To detect spatially localized blur,
each image was divided into a $4\times4$ grid; an image was excluded when more
than 40\% of informative tiles had a tile-level Laplacian variance below 12, or
when the whole-image Laplacian variance was below 48. Tunnel-like views were
identified among non-black images using the conjunction
\begin{equation}
  \mu_Y < 100, \qquad
  \mu_{Y,\mathrm{top}} < 95, \qquad
  V_{\nabla^2Y} < 280,
  \label{eq:tunnel_gate}
\end{equation}
and either a red--blue mean difference greater than 4 or a highlight-glow
contrast greater than 42. The final rejection indicator was the union of the
learned artefact decision, the deterministic rules and image-read failure.

This procedure removed 5,182 of the 384,000 sampled images (1.35\%), leaving
378,818 images for model training and downstream analysis. Because rejection
reasons were not mutually exclusive, their counts were allowed to overlap. The
complete retained and excluded manifests were preserved to make the filtering
decision auditable.

\subsection{Dense DINOv3 feature extraction}
\label{sec:dinov3_features}

We used a frozen DINOv3 ViT-B/16 model pretrained on LVD-1689M as the visual
backbone \cite{simeoni2025dinov3}. Images were converted to RGB, resized to
$224\times224$ pixels with bicubic interpolation and normalized using the
ImageNet channel means $(0.485,0.456,0.406)$ and standard deviations
$(0.229,0.224,0.225)$. Given the patch size of 16 pixels, the image was
represented by a $14\times14$ grid of 196 patch tokens. We removed the class
token and all register tokens from the final hidden sequence and retained the
768-dimensional token associated with each spatial patch. Each token was
independently $\ell_2$-normalized:
\begin{equation}
  \mathbf{t}_{ip}
  = \frac{f_{\theta}(I_i)_p}
  {\lVert f_{\theta}(I_i)_p\rVert_2},
  \qquad
  \mathbf{T}_i\in\mathbb{R}^{196\times768},
  \label{eq:dino_tokens}
\end{equation}
where $i$ indexes images, $p$ indexes spatial patches, and $f_{\theta}$ denotes
the frozen DINOv3 encoder. Patch tensors were stored in half precision using
memory-mapped arrays; 75,264,000 tokens were extracted before quality filtering,
of which 74,248,328 belonged to retained images.

Feature extraction used batches of 128 images and bfloat16 automatic mixed
precision. Inference batches were recursively bisected following a CUDA
out-of-memory exception, providing a deterministic fallback for limited GPU
memory. Extraction was resumable at image level and required approximately
0.95~GiB of peak allocated GPU memory in the recorded run.

\subsection{Feature-level masked autoencoder}
\label{sec:feature_mae}

\subsubsection{Architecture}

We trained a feature-level masked autoencoder on the frozen DINOv3 patch tokens.
The design follows the asymmetric visible-token encoder and lightweight decoder
principle of masked autoencoding \cite{he2022maskedautoencoders}, but reconstructs
DINOv3 features rather than RGB pixels. For image $i$, an affine projection maps
each input token to a 512-dimensional latent space:
\begin{equation}
  \mathbf{z}^{(0)}_{ip}
  = \mathbf{W}_{\mathrm{in}}\mathbf{t}_{ip}
    + \mathbf{b}_{\mathrm{in}} + \boldsymbol{\pi}^{(e)}_p,
  \label{eq:mae_projection}
\end{equation}
where $\boldsymbol{\pi}^{(e)}_p$ is a fixed two-dimensional sine--cosine
positional embedding. For every training image, 75\% of the 196 positions were
randomly masked, leaving 49 visible tokens. Only visible tokens were processed
by the encoder:
\begin{equation}
  \mathbf{H}_{i,V}
  = \operatorname{LN}\!\left(
      E_{\phi}(\mathbf{Z}^{(0)}_{i,V})
    \right),
  \label{eq:mae_encoder}
\end{equation}
where $V$ is the visible-position set. The encoder consisted of four pre-norm
transformer blocks with a hidden width of 512, eight attention heads, a
feed-forward expansion ratio of four, GELU activations and no dropout.

Encoded visible tokens were projected from 512 to 256 dimensions. A learned
mask token was inserted at each missing position, the original spatial order was
restored, and a fixed 256-dimensional sine--cosine positional embedding was
added. A two-layer, eight-head transformer decoder with width 256 and a
fourfold feed-forward expansion produced reconstructed 768-dimensional patch
features:
\begin{equation}
  \widehat{\mathbf{T}}_i
  = \mathbf{W}_{\mathrm{out}}
    \operatorname{LN}\!\left[
      D_{\psi}\!\left(
        \operatorname{Restore}(
          \mathbf{W}_{ed}\mathbf{H}_{i,V},
          \mathbf{m})
        + \boldsymbol{\Pi}^{(d)}
      \right)
    \right]
    + \mathbf{b}_{\mathrm{out}},
  \label{eq:mae_decoder}
\end{equation}
where $\mathbf{m}$ is the learned mask token and
$\mathbf{W}_{ed}$ is the encoder-to-decoder projection. Linear layers were
initialized with Xavier-uniform weights, LayerNorm scales were initialized to
one, and the mask token was sampled from
$\mathcal{N}(0,0.02^2)$. The complete model contained 14,913,280 trainable
parameters.

\subsubsection{Training objective}

Because input patch tokens were normalized, reconstruction was optimized using
cosine distance only at masked positions. If $M_i$ is the masked-position set,
the loss for a batch $\mathcal{B}$ was
\begin{equation}
  \mathcal{L}_{\mathrm{mask}}
  = \frac{1}{\sum_{i\in\mathcal{B}}|M_i|}
    \sum_{i\in\mathcal{B}}
    \sum_{p\in M_i}
    \left[
      1-
      \frac{\mathbf{t}_{ip}^{\top}\widehat{\mathbf{t}}_{ip}}
      {\lVert\mathbf{t}_{ip}\rVert_2
       \lVert\widehat{\mathbf{t}}_{ip}\rVert_2}
    \right].
  \label{eq:masked_cosine_loss}
\end{equation}
No semantic supervision, city label, Top-$K$ operator, non-negativity penalty or
clustering loss was used to train the final Feature-MAE.

\subsubsection{City-balanced optimization}

We reserved 128 clean images per city for validation, producing a validation
set of 3,840 images and a distinct-image training pool of 374,978 images. Each
training batch contained the same number of images from every city. Cities with
fewer retained images were reshuffled and resampled within an epoch until they
matched the largest city-specific training set, preventing cities with more
retained images from dominating optimization.

The model was trained for ten epochs using AdamW with
$\beta_1=0.9$, $\beta_2=0.95$, an initial learning rate of
$1.5\times10^{-4}$ and weight decay 0.05. A linear warm-up covered the first
5\% of optimizer steps, followed by cosine decay to zero. Gradients were clipped
to a global norm of one. Training used bfloat16 automatic mixed precision and
TensorFloat-32 matrix multiplication. The city-balanced batch contained eight
images per city (240 images) during epochs 1--4 and 16 images per city (480
images) during epochs 5--10; the latter was processed as two microbatches of 240
images. The checkpoint with the lowest validation loss was retained. The selected
checkpoint was obtained at epoch 10, with masked cosine losses of 0.10356 and
0.10334 on the training and validation sets, respectively. Peak allocated GPU
memory was 3.27~GiB.

\subsection{Latent activation maps and representative images}
\label{sec:latent_heatmaps}

At inference, masking and the decoder were removed. All 196 DINOv3 patch tokens
were projected and processed by the trained encoder:
\begin{equation}
  \mathbf{A}_i
  = \operatorname{LN}\!\left[
      E_{\phi}(\mathbf{W}_{\mathrm{in}}\mathbf{T}_i
      +\mathbf{b}_{\mathrm{in}}+\boldsymbol{\Pi}^{(e)})
    \right]
  \in \mathbb{R}^{196\times512}.
  \label{eq:full_encoding}
\end{equation}
Reshaping $\mathbf{A}_i$ yielded a
$14\times14\times512$ spatial activation tensor. The activations were signed;
therefore, dimension $d$ was visualized using its positive evidence
$\max(A_{ipd},0)$ rather than the absolute activation. For each dimension, we
defined the image-level response
\begin{equation}
  s_{id}=\max_{p\in\{1,\ldots,196\}} A_{ipd}.
  \label{eq:image_response}
\end{equation}
All 378,818 retained images were scanned. The 128 highest-scoring candidates
were retained provisionally, after which duplicate panoramas were removed and
the ten highest-scoring unique panoramas were selected. This ensured that a
dimension was not interpreted from multiple headings of the same location.

For visualization, positive activations were divided by the maximum response of
the selected examples, clipped to $[0,1]$, mapped to a perceptually ordered
blue--cyan--yellow--red colour scale, and resized to the original image size
using bicubic interpolation. Transparency increased monotonically with
activation so that near-zero patches remained visually unchanged. The resulting
heatmap therefore indicates where a dimension responds, whereas the original
image remains the source for interpreting scene content.

\subsection{Multiview similarity between latent dimensions}
\label{sec:multiview_similarity}

A single parameter-space similarity is insufficient to characterize a contextual
latent dimension. We therefore represented each of the 512 dimensions through
five complementary views: encoder direction, linearized decoder direction,
patch-position profile, city profile and activation-context profile.

First, the encoder direction of dimension $d$ was the $d$th row of the input
projection matrix,
\begin{equation}
  \mathbf{e}_d = (\mathbf{W}_{\mathrm{in}})_{d,:}
  \in\mathbb{R}^{768}.
  \label{eq:encoder_direction}
\end{equation}
Second, we linearized the path from encoder latents to reconstructed DINOv3
features, obtaining
\begin{equation}
  \mathbf{g}_d
  = \left[
      (\mathbf{W}_{\mathrm{out}}\mathbf{W}_{ed})^{\top}
    \right]_{d,:}
  \in\mathbb{R}^{768}.
  \label{eq:decoder_direction}
\end{equation}

For the remaining views, each patch was assigned to its largest latent
activation,
\begin{equation}
  w_{ip}=\underset{d\in\{1,\ldots,512\}}{\arg\max}\; A_{ipd}.
  \label{eq:winning_dimension}
\end{equation}
The patch-position profile
$\mathbf{p}_d\in\mathbb{R}^{196}$ counted how often dimension $d$ won at each
grid position, while the city profile
$\mathbf{c}_d\in\mathbb{R}^{30}$ counted winning patches by city:
\begin{align}
  p_{dq} &= \sum_i \mathbb{I}(w_{iq}=d), \\
  c_{du} &= \sum_{i:\,\operatorname{city}(i)=u}\sum_p
            \mathbb{I}(w_{ip}=d).
  \label{eq:position_city_profiles}
\end{align}

Finally, we constructed an image-level activation-context view. Let
$B_{id}=1$ if dimension $d$ won at least one patch in image $i$, and zero
otherwise. For $N$ images, define the co-occurrence count
$n_{de}=\sum_iB_{id}B_{ie}$ and support $n_d=\sum_iB_{id}$. We used a smoothed
positive pointwise mutual information matrix
\begin{equation}
  Q_{de}=\max\left\{
    0,
    \log\frac{n_{de}+0.5}{n_dn_e/N+0.5}
  \right\}, \qquad Q_{dd}=0,
  \label{eq:ppmi}
\end{equation}
and treated row $\mathbf{q}_d$ as the contextual profile of dimension $d$.

For each view $v\in\{\mathbf{e},\mathbf{g},\mathbf{p},\mathbf{c},\mathbf{q}\}$,
we computed a $512\times512$ cosine-similarity matrix:
\begin{equation}
  S^{(v)}_{de}
  =\frac{\mathbf{v}_d^{\top}\mathbf{v}_e}
  {\lVert\mathbf{v}_d\rVert_2\lVert\mathbf{v}_e\rVert_2}.
  \label{eq:view_cosine}
\end{equation}

\subsection{Rank fusion and hierarchical clustering}
\label{sec:hierarchical_clustering}

The five similarity views have different numerical distributions, so we fused
neighbour ranks rather than raw similarity values. For every dimension $d$ and
view $v$, the 30 nearest non-self neighbours were ordered by decreasing
similarity. A neighbour at zero-indexed rank $r\in\{0,\ldots,29\}$ received
weight $1-r/30$. The directed fused affinity was
\begin{equation}
  \widetilde{F}_{de}
  =\frac{1}{5}\sum_v
  \mathbb{I}\!\left[e\in\mathcal{N}^{(v)}_{30}(d)\right]
  \left(1-\frac{r^{(v)}_{de}}{30}\right).
  \label{eq:rank_fusion}
\end{equation}
We then symmetrized and rescaled the matrix:
\begin{equation}
  F=\frac{\widetilde{F}+\widetilde{F}^{\top}}{2},
  \qquad F_{dd}=0, \qquad
  F\leftarrow\frac{F}{\max_{d,e}F_{de}}.
  \label{eq:fused_affinity}
\end{equation}

To preserve nonlinear neighbourhood structure, the fused affinity was projected
into 32 dimensions using the eigenvectors of the symmetric normalized graph
Laplacian \cite{ng2001spectral}. The first eigenvector was retained, consistent
with the spectral-clustering embedding used in the implementation. Ward's
minimum-variance agglomerative method was then applied to Euclidean distances in
the spectral embedding \cite{ward1963hierarchical}. One dendrogram was
constructed and cut at two resolutions: 32 coarse clusters and 64 fine clusters.
Because both partitions were cuts of the same tree, each fine cluster had
exactly one coarse parent by construction.

For cluster $C$, the representative or medoid dimension was defined using the
fused affinity rather than the spectral coordinates:
\begin{equation}
  d_C^{*}=\underset{d\in C}{\arg\max}
  \sum_{e\in C}F_{de}.
  \label{eq:cluster_medoid}
\end{equation}
Cluster identities were assigned only after the tree had been fixed. Canonical
identifiers C000--C031 and F000--F063 were ordered by the first latent dimension
appearing in each cluster and carry no semantic meaning.

\subsection{Cluster validity and split-sample stability}
\label{sec:cluster_validation}

We characterized both hierarchy levels using silhouette score,
Calinski--Harabasz index, Davies--Bouldin index and normalized cluster-size
entropy. If $n_k$ denotes the size of cluster $k$, $p_k=n_k/512$ and $K$ is the
number of clusters, normalized size entropy was
\begin{equation}
  H_{\mathrm{size}}
  =-\frac{\sum_{k=1}^{K}p_k\log p_k}{\log K}.
  \label{eq:size_entropy}
\end{equation}

Stability was evaluated using two deterministic, interleaved halves of the image
sample. Parameter-space views (encoder and decoder directions) were held fixed,
whereas position, city and PPMI-context profiles were independently recomputed
for each half. The complete rank-fusion, spectral-embedding and Ward-clustering
pipeline was repeated for each half using seeds 43 and 44. For each dimension,
we calculated the Jaccard overlap between its full-data cluster and its
corresponding half-sample cluster. Global agreement was additionally summarized
with the adjusted Rand index and normalized mutual information at both the
32- and 64-cluster resolutions. These stability measures quantify sensitivity to
the image sample and were not used to choose or modify semantic labels.

\subsection{Post-hoc semantic interpretation}
\label{sec:semantic_interpretation}

We used the locally deployed Qwen3-VL-2B-Instruct model
\cite{bai2025qwen3vl} to attach human-readable descriptions to the statistical
clusters. Evidence sheets were resized to a maximum side length of 1,200 pixels,
and generation was limited to 240 new tokens. Interpretation was performed after
Feature-MAE training and after the
32/64 hierarchy had been frozen. The language model therefore had no access to
the optimization objective and could not influence dimension similarity or
cluster membership.

For each of the 64 fine and 32 coarse clusters, we selected the cluster medoid
from Eq.~\ref{eq:cluster_medoid}. We then retrieved that dimension's ten
highest-scoring images from unique panoramas. Each item in the evidence sheet
contained (i) the original street-view image, (ii) the activation overlay and
(iii) a crop centred on the maximally activated patch. The model was instructed
to identify visual evidence that recurred across the ten examples and to return
a strict JSON object containing a concise Chinese name, a short English label,
a semantic type, a summary, three visual cues, a tentative urban interpretation,
cross-image consistency, artefact probability and confidence. Semantic type was
restricted to object, material, colour/light, texture, geometry, spatial layout,
environment, artefact, mixed or unclear.

Fine clusters were interpreted first. When a coarse cluster was interpreted,
the already generated descriptions of its fine children were supplied together
with the coarse medoid evidence, allowing the model to formulate a broader
parent concept without changing the statistical hierarchy. Outputs were parsed
as JSON and automatically retried when parsing failed. A second generation was
requested when a name contained prohibited generic terms, confidence was
mechanically set near zero, or the reported visual cues were duplicated. Exactly
ten images were used for each of the 96 categories; interpreting every training
image was neither necessary nor performed.

\subsection{Implementation and reproducibility}
\label{sec:implementation}

All stochastic components used a global seed of 42 unless a split-sample seed
was explicitly specified. DINOv3 and Qwen3-VL weights were loaded locally, and
the DINOv3 backbone remained frozen throughout. Large arrays were stored as
half-precision memory maps, and image selection, feature extraction, training,
ranking and clustering wrote resumable intermediate states. GPU-intensive
PyTorch stages and CPU-intensive SciPy/scikit-learn stages were executed in
separate processes to avoid conflicts between OpenMP runtimes. The statistical
hierarchy, category membership table, linkage matrix, fused similarity matrix,
spectral embedding, heatmap metadata, semantic prompts and raw language-model
responses were retained to support independent auditing and figure
reconstruction.


\section{Methodology}
\label{sec:methodology}

\subsection{Overview}
\label{sec:method_overview}

We developed a fully unsupervised framework for discovering multiscale urban
visual patterns from geotagged street-view imagery. The framework contains six
stages: (i) geographically broad and city-balanced image sampling; (ii) image
quality control; (iii) extraction of dense patch representations using a frozen
DINOv3 vision transformer; (iv) contextual representation learning using a
feature-level masked autoencoder (Feature-MAE); (v) spatial visualization and
hierarchical organization of the learned latent dimensions; and (vi) post-hoc
semantic interpretation using a compact vision--language model. Statistical
learning and clustering were completed before semantic labels were generated,
so language-model outputs could not affect the learned representation or the
cluster assignments.

The final analysis used Feature-MAE rather than the Top-$K$ sparse autoencoder
examined during model development. Both models used a latent width of 512, but
the preliminary sparse autoencoder retained the 32 largest activations for each
input, whereas the final Feature-MAE learned signed, contextual patch codes by
reconstructing masked DINOv3 features. Unless explicitly described as a pilot
or ablation, all analyses and figures reported in this study refer to the final
Feature-MAE model.

\subsection{Street-view image sample}
\label{sec:streetview_sample}

The study sample comprised street-view imagery from 30 cities spanning Africa,
Asia, Europe, North America, South America and Oceania: Amsterdam, Bangkok,
Bogot\'a, Buenos Aires, Cape Town, Dhaka, Hong Kong, Istanbul, Jakarta,
Johannesburg, Kuala Lumpur, Lagos, Lima, London, Los Angeles, Manila, Mexico
City, Moscow, Mumbai, New Delhi, New York City, Osaka, Paris, S\~ao Paulo,
Seoul, Singapore, Sydney, Taipei, Toronto and Vienna. Four nominal viewing
directions ($0^{\circ}$, $90^{\circ}$, $180^{\circ}$ and $270^{\circ}$) were
available for each panorama.

We first sampled 10,000 panoramas per city. From their directional views, we
then selected 12,800 images per city, yielding 384,000 images before quality
control. Selection used a deterministic, city-specific random seed derived from
the global seed of 42. Within each city, panorama order and heading order were
randomized and images were selected in a round-robin procedure: one randomly
ordered heading was selected from every panorama before a second heading was
selected from any panorama. This design retained all 10,000 unique panoramas in
each city while avoiding disproportionate sampling of multiple headings from a
small number of locations. It also maintained approximately balanced heading
counts within each city.

\subsection{Image quality control}
\label{sec:image_qc}

Image quality was assessed before representation learning to reduce the risk
that latent dimensions would encode acquisition defects rather than urban
content. Each image was resized to $256\times256$ pixels and described using 11
low-cost photometric and structural measurements: mean and standard deviation
of luminance, saturated-pixel fraction, dark-pixel fraction, Laplacian variance,
edge density, luminance entropy, dominant-intensity fraction, colourfulness,
mean red--green contrast and mean yellow--blue contrast. A histogram-based
gradient-boosting classifier, trained from conservative weak labels for clear
good- and bad-quality images, converted these measurements into an artefact
probability $p_{\mathrm{bad}}$. Images with $p_{\mathrm{bad}}>0.5$ were excluded.

We supplemented this classifier with explicit rules for unreadable, black,
blurred and tunnel-like frames. Let $\mu_Y$, $\sigma_Y$, $f_{Y<20}$ and
$V_{\nabla^2Y}$ denote mean luminance, luminance standard deviation, the
fraction of pixels with luminance below 20, and Laplacian variance, respectively.
An image was classified as black or uninformative if
$\mu_Y<14$, $f_{Y<20}>0.97$ or $\sigma_Y<6$. To detect spatially localized blur,
each image was divided into a $4\times4$ grid; an image was excluded when more
than 40\% of informative tiles had a tile-level Laplacian variance below 12, or
when the whole-image Laplacian variance was below 48. Tunnel-like views were
identified among non-black images using the conjunction
\begin{equation}
  \mu_Y < 100, \qquad
  \mu_{Y,\mathrm{top}} < 95, \qquad
  V_{\nabla^2Y} < 280,
  \label{eq:tunnel_gate}
\end{equation}
and either a red--blue mean difference greater than 4 or a highlight-glow
contrast greater than 42. The final rejection indicator was the union of the
learned artefact decision, the deterministic rules and image-read failure.

This procedure removed 5,182 of the 384,000 sampled images (1.35\%), leaving
378,818 images for model training and downstream analysis. Because rejection
reasons were not mutually exclusive, their counts were allowed to overlap. The
complete retained and excluded manifests were preserved to make the filtering
decision auditable.

\subsection{Dense DINOv3 feature extraction}
\label{sec:dinov3_features}

We used a frozen DINOv3 ViT-B/16 model pretrained on LVD-1689M as the visual
backbone \cite{simeoni2025dinov3}. Images were converted to RGB, resized to
$224\times224$ pixels with bicubic interpolation and normalized using the
ImageNet channel means $(0.485,0.456,0.406)$ and standard deviations
$(0.229,0.224,0.225)$. Given the patch size of 16 pixels, the image was
represented by a $14\times14$ grid of 196 patch tokens. We removed the class
token and all register tokens from the final hidden sequence and retained the
768-dimensional token associated with each spatial patch. Each token was
independently $\ell_2$-normalized:
\begin{equation}
  \mathbf{t}_{ip}
  = \frac{f_{\theta}(I_i)_p}
  {\lVert f_{\theta}(I_i)_p\rVert_2},
  \qquad
  \mathbf{T}_i\in\mathbb{R}^{196\times768},
  \label{eq:dino_tokens}
\end{equation}
where $i$ indexes images, $p$ indexes spatial patches, and $f_{\theta}$ denotes
the frozen DINOv3 encoder. Patch tensors were stored in half precision using
memory-mapped arrays; 75,264,000 tokens were extracted before quality filtering,
of which 74,248,328 belonged to retained images.

Feature extraction used batches of 128 images and bfloat16 automatic mixed
precision. Inference batches were recursively bisected following a CUDA
out-of-memory exception, providing a deterministic fallback for limited GPU
memory. Extraction was resumable at image level and required approximately
0.95~GiB of peak allocated GPU memory in the recorded run.

\subsection{Feature-level masked autoencoder}
\label{sec:feature_mae}

\subsubsection{Architecture}

We trained a feature-level masked autoencoder on the frozen DINOv3 patch tokens.
The design follows the asymmetric visible-token encoder and lightweight decoder
principle of masked autoencoding \cite{he2022maskedautoencoders}, but reconstructs
DINOv3 features rather than RGB pixels. For image $i$, an affine projection maps
each input token to a 512-dimensional latent space:
\begin{equation}
  \mathbf{z}^{(0)}_{ip}
  = \mathbf{W}_{\mathrm{in}}\mathbf{t}_{ip}
    + \mathbf{b}_{\mathrm{in}} + \boldsymbol{\pi}^{(e)}_p,
  \label{eq:mae_projection}
\end{equation}
where $\boldsymbol{\pi}^{(e)}_p$ is a fixed two-dimensional sine--cosine
positional embedding. For every training image, 75\% of the 196 positions were
randomly masked, leaving 49 visible tokens. Only visible tokens were processed
by the encoder:
\begin{equation}
  \mathbf{H}_{i,V}
  = \operatorname{LN}\!\left(
      E_{\phi}(\mathbf{Z}^{(0)}_{i,V})
    \right),
  \label{eq:mae_encoder}
\end{equation}
where $V$ is the visible-position set. The encoder consisted of four pre-norm
transformer blocks with a hidden width of 512, eight attention heads, a
feed-forward expansion ratio of four, GELU activations and no dropout.

Encoded visible tokens were projected from 512 to 256 dimensions. A learned
mask token was inserted at each missing position, the original spatial order was
restored, and a fixed 256-dimensional sine--cosine positional embedding was
added. A two-layer, eight-head transformer decoder with width 256 and a
fourfold feed-forward expansion produced reconstructed 768-dimensional patch
features:
\begin{equation}
  \widehat{\mathbf{T}}_i
  = \mathbf{W}_{\mathrm{out}}
    \operatorname{LN}\!\left[
      D_{\psi}\!\left(
        \operatorname{Restore}(
          \mathbf{W}_{ed}\mathbf{H}_{i,V},
          \mathbf{m})
        + \boldsymbol{\Pi}^{(d)}
      \right)
    \right]
    + \mathbf{b}_{\mathrm{out}},
  \label{eq:mae_decoder}
\end{equation}
where $\mathbf{m}$ is the learned mask token and
$\mathbf{W}_{ed}$ is the encoder-to-decoder projection. Linear layers were
initialized with Xavier-uniform weights, LayerNorm scales were initialized to
one, and the mask token was sampled from
$\mathcal{N}(0,0.02^2)$. The complete model contained 14,913,280 trainable
parameters.

\subsubsection{Training objective}

Because input patch tokens were normalized, reconstruction was optimized using
cosine distance only at masked positions. If $M_i$ is the masked-position set,
the loss for a batch $\mathcal{B}$ was
\begin{equation}
  \mathcal{L}_{\mathrm{mask}}
  = \frac{1}{\sum_{i\in\mathcal{B}}|M_i|}
    \sum_{i\in\mathcal{B}}
    \sum_{p\in M_i}
    \left[
      1-
      \frac{\mathbf{t}_{ip}^{\top}\widehat{\mathbf{t}}_{ip}}
      {\lVert\mathbf{t}_{ip}\rVert_2
       \lVert\widehat{\mathbf{t}}_{ip}\rVert_2}
    \right].
  \label{eq:masked_cosine_loss}
\end{equation}
No semantic supervision, city label, Top-$K$ operator, non-negativity penalty or
clustering loss was used to train the final Feature-MAE.

\subsubsection{City-balanced optimization}

We reserved 128 clean images per city for validation, producing a validation
set of 3,840 images and a distinct-image training pool of 374,978 images. Each
training batch contained the same number of images from every city. Cities with
fewer retained images were reshuffled and resampled within an epoch until they
matched the largest city-specific training set, preventing cities with more
retained images from dominating optimization.

The model was trained for ten epochs using AdamW with
$\beta_1=0.9$, $\beta_2=0.95$, an initial learning rate of
$1.5\times10^{-4}$ and weight decay 0.05. A linear warm-up covered the first
5\% of optimizer steps, followed by cosine decay to zero. Gradients were clipped
to a global norm of one. Training used bfloat16 automatic mixed precision and
TensorFloat-32 matrix multiplication. The city-balanced batch contained eight
images per city (240 images) during epochs 1--4 and 16 images per city (480
images) during epochs 5--10; the latter was processed as two microbatches of 240
images. The checkpoint with the lowest validation loss was retained. The selected
checkpoint was obtained at epoch 10, with masked cosine losses of 0.10356 and
0.10334 on the training and validation sets, respectively. Peak allocated GPU
memory was 3.27~GiB.

\subsection{Latent activation maps and representative images}
\label{sec:latent_heatmaps}

At inference, masking and the decoder were removed. All 196 DINOv3 patch tokens
were projected and processed by the trained encoder:
\begin{equation}
  \mathbf{A}_i
  = \operatorname{LN}\!\left[
      E_{\phi}(\mathbf{W}_{\mathrm{in}}\mathbf{T}_i
      +\mathbf{b}_{\mathrm{in}}+\boldsymbol{\Pi}^{(e)})
    \right]
  \in \mathbb{R}^{196\times512}.
  \label{eq:full_encoding}
\end{equation}
Reshaping $\mathbf{A}_i$ yielded a
$14\times14\times512$ spatial activation tensor. The activations were signed;
therefore, dimension $d$ was visualized using its positive evidence
$\max(A_{ipd},0)$ rather than the absolute activation. For each dimension, we
defined the image-level response
\begin{equation}
  s_{id}=\max_{p\in\{1,\ldots,196\}} A_{ipd}.
  \label{eq:image_response}
\end{equation}
All 378,818 retained images were scanned. The 128 highest-scoring candidates
were retained provisionally, after which duplicate panoramas were removed and
the ten highest-scoring unique panoramas were selected. This ensured that a
dimension was not interpreted from multiple headings of the same location.

For visualization, positive activations were divided by the maximum response of
the selected examples, clipped to $[0,1]$, mapped to a perceptually ordered
blue--cyan--yellow--red colour scale, and resized to the original image size
using bicubic interpolation. Transparency increased monotonically with
activation so that near-zero patches remained visually unchanged. The resulting
heatmap therefore indicates where a dimension responds, whereas the original
image remains the source for interpreting scene content.

\subsection{Multiview similarity between latent dimensions}
\label{sec:multiview_similarity}

A single parameter-space similarity is insufficient to characterize a contextual
latent dimension. We therefore represented each of the 512 dimensions through
five complementary views: encoder direction, linearized decoder direction,
patch-position profile, city profile and activation-context profile.

First, the encoder direction of dimension $d$ was the $d$th row of the input
projection matrix,
\begin{equation}
  \mathbf{e}_d = (\mathbf{W}_{\mathrm{in}})_{d,:}
  \in\mathbb{R}^{768}.
  \label{eq:encoder_direction}
\end{equation}
Second, we linearized the path from encoder latents to reconstructed DINOv3
features, obtaining
\begin{equation}
  \mathbf{g}_d
  = \left[
      (\mathbf{W}_{\mathrm{out}}\mathbf{W}_{ed})^{\top}
    \right]_{d,:}
  \in\mathbb{R}^{768}.
  \label{eq:decoder_direction}
\end{equation}

For the remaining views, each patch was assigned to its largest latent
activation,
\begin{equation}
  w_{ip}=\underset{d\in\{1,\ldots,512\}}{\arg\max}\; A_{ipd}.
  \label{eq:winning_dimension}
\end{equation}
The patch-position profile
$\mathbf{p}_d\in\mathbb{R}^{196}$ counted how often dimension $d$ won at each
grid position, while the city profile
$\mathbf{c}_d\in\mathbb{R}^{30}$ counted winning patches by city:
\begin{align}
  p_{dq} &= \sum_i \mathbb{I}(w_{iq}=d), \\
  c_{du} &= \sum_{i:\,\operatorname{city}(i)=u}\sum_p
            \mathbb{I}(w_{ip}=d).
  \label{eq:position_city_profiles}
\end{align}

Finally, we constructed an image-level activation-context view. Let
$B_{id}=1$ if dimension $d$ won at least one patch in image $i$, and zero
otherwise. For $N$ images, define the co-occurrence count
$n_{de}=\sum_iB_{id}B_{ie}$ and support $n_d=\sum_iB_{id}$. We used a smoothed
positive pointwise mutual information matrix
\begin{equation}
  Q_{de}=\max\left\{
    0,
    \log\frac{n_{de}+0.5}{n_dn_e/N+0.5}
  \right\}, \qquad Q_{dd}=0,
  \label{eq:ppmi}
\end{equation}
and treated row $\mathbf{q}_d$ as the contextual profile of dimension $d$.

For each view $v\in\{\mathbf{e},\mathbf{g},\mathbf{p},\mathbf{c},\mathbf{q}\}$,
we computed a $512\times512$ cosine-similarity matrix:
\begin{equation}
  S^{(v)}_{de}
  =\frac{\mathbf{v}_d^{\top}\mathbf{v}_e}
  {\lVert\mathbf{v}_d\rVert_2\lVert\mathbf{v}_e\rVert_2}.
  \label{eq:view_cosine}
\end{equation}

\subsection{Rank fusion and hierarchical clustering}
\label{sec:hierarchical_clustering}

The five similarity views have different numerical distributions, so we fused
neighbour ranks rather than raw similarity values. For every dimension $d$ and
view $v$, the 30 nearest non-self neighbours were ordered by decreasing
similarity. A neighbour at zero-indexed rank $r\in\{0,\ldots,29\}$ received
weight $1-r/30$. The directed fused affinity was
\begin{equation}
  \widetilde{F}_{de}
  =\frac{1}{5}\sum_v
  \mathbb{I}\!\left[e\in\mathcal{N}^{(v)}_{30}(d)\right]
  \left(1-\frac{r^{(v)}_{de}}{30}\right).
  \label{eq:rank_fusion}
\end{equation}
We then symmetrized and rescaled the matrix:
\begin{equation}
  F=\frac{\widetilde{F}+\widetilde{F}^{\top}}{2},
  \qquad F_{dd}=0, \qquad
  F\leftarrow\frac{F}{\max_{d,e}F_{de}}.
  \label{eq:fused_affinity}
\end{equation}

To preserve nonlinear neighbourhood structure, the fused affinity was projected
into 32 dimensions using the eigenvectors of the symmetric normalized graph
Laplacian \cite{ng2001spectral}. The first eigenvector was retained, consistent
with the spectral-clustering embedding used in the implementation. Ward's
minimum-variance agglomerative method was then applied to Euclidean distances in
the spectral embedding \cite{ward1963hierarchical}. One dendrogram was
constructed and cut at two resolutions: 32 coarse clusters and 64 fine clusters.
Because both partitions were cuts of the same tree, each fine cluster had
exactly one coarse parent by construction.

For cluster $C$, the representative or medoid dimension was defined using the
fused affinity rather than the spectral coordinates:
\begin{equation}
  d_C^{*}=\underset{d\in C}{\arg\max}
  \sum_{e\in C}F_{de}.
  \label{eq:cluster_medoid}
\end{equation}
Cluster identities were assigned only after the tree had been fixed. Canonical
identifiers C000--C031 and F000--F063 were ordered by the first latent dimension
appearing in each cluster and carry no semantic meaning.

\subsection{Cluster validity and split-sample stability}
\label{sec:cluster_validation}

We characterized both hierarchy levels using silhouette score,
Calinski--Harabasz index, Davies--Bouldin index and normalized cluster-size
entropy. If $n_k$ denotes the size of cluster $k$, $p_k=n_k/512$ and $K$ is the
number of clusters, normalized size entropy was
\begin{equation}
  H_{\mathrm{size}}
  =-\frac{\sum_{k=1}^{K}p_k\log p_k}{\log K}.
  \label{eq:size_entropy}
\end{equation}

Stability was evaluated using two deterministic, interleaved halves of the image
sample. Parameter-space views (encoder and decoder directions) were held fixed,
whereas position, city and PPMI-context profiles were independently recomputed
for each half. The complete rank-fusion, spectral-embedding and Ward-clustering
pipeline was repeated for each half using seeds 43 and 44. For each dimension,
we calculated the Jaccard overlap between its full-data cluster and its
corresponding half-sample cluster. Global agreement was additionally summarized
with the adjusted Rand index and normalized mutual information at both the
32- and 64-cluster resolutions. These stability measures quantify sensitivity to
the image sample and were not used to choose or modify semantic labels.

\subsection{Post-hoc semantic interpretation}
\label{sec:semantic_interpretation}

We used the locally deployed Qwen3-VL-2B-Instruct model
\cite{bai2025qwen3vl} to attach human-readable descriptions to the statistical
clusters. Evidence sheets were resized to a maximum side length of 1,200 pixels,
and generation was limited to 240 new tokens. Interpretation was performed after
Feature-MAE training and after the
32/64 hierarchy had been frozen. The language model therefore had no access to
the optimization objective and could not influence dimension similarity or
cluster membership.

For each of the 64 fine and 32 coarse clusters, we selected the cluster medoid
from Eq.~\ref{eq:cluster_medoid}. We then retrieved that dimension's ten
highest-scoring images from unique panoramas. Each item in the evidence sheet
contained (i) the original street-view image, (ii) the activation overlay and
(iii) a crop centred on the maximally activated patch. The model was instructed
to identify visual evidence that recurred across the ten examples and to return
a strict JSON object containing a concise Chinese name, a short English label,
a semantic type, a summary, three visual cues, a tentative urban interpretation,
cross-image consistency, artefact probability and confidence. Semantic type was
restricted to object, material, colour/light, texture, geometry, spatial layout,
environment, artefact, mixed or unclear.

Fine clusters were interpreted first. When a coarse cluster was interpreted,
the already generated descriptions of its fine children were supplied together
with the coarse medoid evidence, allowing the model to formulate a broader
parent concept without changing the statistical hierarchy. Outputs were parsed
as JSON and automatically retried when parsing failed. A second generation was
requested when a name contained prohibited generic terms, confidence was
mechanically set near zero, or the reported visual cues were duplicated. Exactly
ten images were used for each of the 96 categories; interpreting every training
image was neither necessary nor performed.

\subsection{Implementation and reproducibility}
\label{sec:implementation}

All stochastic components used a global seed of 42 unless a split-sample seed
was explicitly specified. DINOv3 and Qwen3-VL weights were loaded locally, and
the DINOv3 backbone remained frozen throughout. Large arrays were stored as
half-precision memory maps, and image selection, feature extraction, training,
ranking and clustering wrote resumable intermediate states. GPU-intensive
PyTorch stages and CPU-intensive SciPy/scikit-learn stages were executed in
separate processes to avoid conflicts between OpenMP runtimes. The statistical
hierarchy, category membership table, linkage matrix, fused similarity matrix,
spectral embedding, heatmap metadata, semantic prompts and raw language-model
responses were retained to support independent auditing and figure
reconstruction.
